\section*{Chapitre \ref{ch:fm:intellectual-property-and-non-competes} --- Propriété intellectuelle et clauses de non-concurrence}

\begin{solution}{exo:fm:intellectual-property-and-non-competes:1}
Avec $\lambda=\ln2$ et $r=0$~: $\text{protégé}(T)=(5/\ln2)(1-2^{-T})$, soit 2.11 pour six mois et 3.61 pour un an~; la perte sans \omterm{def:fm:intellectual-property-and-non-competes:restrict}{garden leave} est
$5/\ln2=7.21$.
\end{solution}

\begin{solution}{exo:fm:intellectual-property-and-non-competes:2}
Non généralement connue ni aisément vérifiable~; valeur commerciale du fait de son caractère secret~; mesures raisonnables pour le garder secret. Le signal publié n'est pas un
secret~; la version calibrée par la société, ses paramètres et sa mise en œuvre peuvent l'être, si la société les garde secrets.
\end{solution}

\begin{solution}{exo:fm:intellectual-property-and-non-competes:3}
$1.5\,(1-e^{-0.08T})/0.08$~: 0.37, 0.74 et 1.44 (millions de dollars).
\end{solution}

\begin{solution}{exo:fm:intellectual-property-and-non-competes:4}
$T^*=h\log_2(LP/S)=6$ mois donnent $h=6/3=2$ mois lorsque $LP/S=8$ et $h=6/2=3$ mois lorsque $LP/S=4$. Six mois est la bonne durée seulement pour
des connaissances dont la valeur décroît en quelques mois.
\end{solution}

\begin{solution}{exo:fm:intellectual-property-and-non-competes:5}
Le tribunal a jugé le code source immatériel, donc pas un «~bien~» volé au sens du National Stolen Property Act, et le système de trading, qui n'a pas été
vendu ni concédé sous licence, donc pas un produit «~fabriqué pour ou placé dans~» le commerce interétatique au sens de l'Economic Espionage Act. La loi de 2012 a remplacé ces
termes par un produit ou un service «~utilisé dans ou destiné à être utilisé dans~» le commerce.
\end{solution}

\begin{solution}{exo:fm:intellectual-property-and-non-competes:6}
$\text{protégé}(T^*)=(LP/\lambda)(1-S/LP)=(LP-S)/\lambda$ et $\text{coût}(T^*)=S\ln(LP/S)/\lambda$, donc la valeur nette vaut $(LP-S-S\ln(LP/S))/\lambda$.
Avec $LP=5$, $S=1$, $\lambda=\ln2$~: $(4-\ln5)/\ln2=3.45$, comme dans l'exemple.
\end{solution}

\begin{solution}{exo:fm:intellectual-property-and-non-competes:7}
Demi-vies de 23.9, 6.6 et 11.7 mois~; durées optimales de 3.99, 1.65 et 1.38 ans. Le classement des trois types est inchangé~; les durées varient de
quelques mois, autant que le bruit dans les demi-vies calibrées.
\end{solution}

\begin{solution}{exo:fm:intellectual-property-and-non-competes:8}
La protection ne dure que le temps de la restriction~: avec une demi-vie de 24.9 mois, deux ans protègent 56~\% de la valeur de la structure pour un concurrent
(valeur nette de 5.38 millions de dollars après le coût du \omterm{def:fm:intellectual-property-and-non-competes:restrict}{garden leave}), et le reste fuit à son terme. Une restriction de deux ans peut aussi ne pas être exécutoire, et elle ne
protège contre aucune fuite passant par d'autres personnes. La structure est protégée surtout en la conservant comme \omterm{def:fm:intellectual-property-and-non-competes:secret}{secret des affaires}.
\end{solution}

\begin{solution}{pb:fm:intellectual-property-and-non-competes:1}
\begin{enumerate}
\item Code source, structure et paramètres du modèle, données nettoyées, journaux de recherche et savoir-faire de ceux qui les ont construits.
\item Une stratégie est rarement brevetable, et un brevet la publierait~; le droit d'auteur protège le texte du code, non l'idée qu'il met en œuvre.
\item Voir la \cref{def:fm:intellectual-property-and-non-competes:secret}~; les mesures raisonnables pour le garder secret.
\item Voir la \cref{def:fm:intellectual-property-and-non-competes:nda}.
\item Un programmeur a téléversé plus de 500\,000 lignes du code du système de trading de la banque le jour de son départ~; il a été reconnu coupable et condamné à 97 mois~;
la cour d'appel du Second Circuit a annulé en 2012 au motif qu'aucun des deux textes, dans leur rédaction, ne couvrait ce comportement.
\item Elle a étendu l'infraction pénale aux secrets liés aux produits ou services utilisés dans le commerce~; la loi de 2016 a donné aux titulaires un droit
d'action civile fédérale.
\item La règle de la FTC, annoncée en avril 2024, bloquée par un tribunal en août 2024~; la FTC a abandonné ses recours en septembre 2025. Elle n'est pas en vigueur.
\item Rupture de contrat, détournement de \omterm{def:fm:intellectual-property-and-non-competes:secret}{secrets des affaires} et demandes connexes portant sur une stratégie de trading emportée par deux employés passés chez Millennium~; cette
procédure a été rejetée avec préjudice par accord des parties en décembre 2024, sans constat sur le fond.
\item Voir la \cref{def:fm:intellectual-property-and-non-competes:restrict}.
\item $LP(1-e^{-(\lambda+r)T})/(\lambda+r)$ et $S(1-e^{-rT})/r$.
\item Voir la \cref{prop:fm:intellectual-property-and-non-competes:six}.
\item 24.9, 5.7 et 12.1 mois (vraies valeurs 24, 6 et 12).
\item 2.71, 4.18 et 1.66 (millions de dollars)~: 19~\%, 54~\% et 32~\% de la perte sans \omterm{def:fm:intellectual-property-and-non-competes:restrict}{garden leave} (14.47, 7.78 et 5.22).
\item 4.14, 1.42 et 1.43 ans~: $\log_2(LP/S)$ demi-vies, soit 2, 3 et 1.42.
\item Pour le modèle, 1.4 ans pour les paramètres et quatre ans pour la structure~; six mois protègent 54~\% de la valeur des paramètres et 19~\% de la
structure, pour un coût de 0.74.
\item Droit du \omterm{def:fm:intellectual-property-and-non-competes:secret}{secret des affaires}, et les mesures raisonnables qui le gardent secret~: contrôle des accès, journalisation et contrats.
\item Couper les accès, relire les journaux des copies et des téléversements, confirmer les obligations par écrit, faire courir le délai de préavis, prendre conseil auprès d'un juriste.
\item Modérément~: un autre tirage du bruit les déplace de quelques mois (3.99, 1.65 et 1.38 ans), mais pas le classement.
\item $T^*=\ln(LP/S)/\lambda=h\log_2(LP/S)$, quel que soit le taux d'actualisation~: autant de demi-vies qu'il en faut pour que les profits en jeu soient divisés par deux jusqu'au
salaire.
\item Garder le secret par le contrôle des accès, la journalisation et les contrats~; fixer le préavis rémunéré et les restrictions selon la vitesse de décroissance de chaque type de connaissance, dans les limites de ce que
la législation locale le permet, et ne s'appuyer sur elles que pour la partie à décroissance rapide.
\end{enumerate}
\end{solution}

\begin{solution}{iq:fm:intellectual-property-and-non-competes:1}
Préavis rémunéré pendant lequel on reste salarié mais tenu à l'écart du travail et des marchés~; il empêche ce que l'on sait d'atteindre un concurrent tant que
l'information est encore fraîche, et coûte moins cher que l'edge qu'il protège.
\iqlookfor{la valeur décroissante de l'information.}
\end{solution}

\begin{solution}{iq:fm:intellectual-property-and-non-competes:2}
Non~: selon une \omterm{def:fm:intellectual-property-and-non-competes:nda}{clause de cession des inventions}, il appartient à l'ancien employeur, et il peut être un \omterm{def:fm:intellectual-property-and-non-competes:secret}{secret des affaires}. On apporte sa compétence générale et
ses connaissances, non le code.
\iqlookfor{propriété, et frontière entre compétence et bien de l'employeur.}
\end{solution}

\begin{solution}{iq:fm:intellectual-property-and-non-competes:3}
La moitié de sa valeur restante, avant actualisation~: $1-2^{-T/h}$ avec $T=h$.
\iqlookfor{$1-e^{-\lambda T}$.}
\end{solution}

\begin{solution}{iq:fm:intellectual-property-and-non-competes:4}
Attributions d'accès et leurs revues, copies et téléversements hors des systèmes de la société, supports amovibles, clones de dépôts, et les
reconnaissances écrites des obligations de confidentialité.
\iqlookfor{preuves de mesures raisonnables.}
\end{solution}

\begin{solution}{iq:fm:intellectual-property-and-non-competes:5}
Les paramètres~: ils sont recalibrés tous les quelques mois. Une restriction courte les protège bien~; la structure exige le secret, car aucune restriction exécutoire
ne dure assez longtemps.
\iqlookfor{demi-vies selon le type de connaissance.}
\end{solution}

\begin{solution}{iq:fm:intellectual-property-and-non-competes:6}
Valeur nette $LP(1-e^{-(\lambda+r)T})/(\lambda+r)-S(1-e^{-rT})/r$~; sa dérivée $e^{-rT}(LPe^{-\lambda T}-S)$ s'annule en $T^*=\ln(LP/S)/\lambda$.
\iqlookfor{protection marginale face au coût marginal, et que $r$ disparaît.}
\end{solution}
